\documentclass[11pt]{article}
\usepackage{amsmath,amssymb,amsthm}
\newtheorem{theorem}{Theorem}
\title{One standby relay, several possible recipients}
\begin{document}
\maketitle

A relay must keep an existing geometric link while it moves to serve
whichever recipient a report names. Each pair of possible recipients may
admit a deadline even though no policy meets it for all recipients. The
distinction is one shared standby state, rather than a new placement for
each pair.

\paragraph{Model.}
Let \(C\subset\mathbb R^d\) be nonempty, compact and convex. A finite
nonempty set of worlds \(J\) has nonempty closed service regions
\(F_j\subseteq C\). A truthful report names \(j\) at time zero.
Before that report there is no observation that distinguishes the worlds.
The relay may have any common position and velocity achievable while
remaining in \(C\); in particular, any placement at rest in \(C\) is
available before the clock starts. From the report to the deadline \(T\),
it must remain in \(C\), satisfy \(\|\dot v\|\le a\) almost everywhere and
\(\|v\|\le V\), and end at rest in \(F_j\). The bounds \(a,V\) are positive
and isotropic; velocity is absolutely continuous. A common stationary
delivery dwell \(w\) after arrival can be budgeted separately, provided
service remains available during that dwell. Obstacles, moving recipients,
link intermittence, and faults during forwarding are not part of the
model. Region membership is not a packet-delivery guarantee.

\begin{theorem}
Define
\[
 r=\min_{x\in C}\max_{j\in J}\operatorname{dist}(x,F_j),\qquad
 D(T)=
 \begin{cases}
 aT^2/4, & T\le 2V/a,\\
 VT-V^2/a, & T\ge 2V/a .
 \end{cases}
\]
There is a causal policy meeting the deadline for every world if and only
if \(r\le D(T)\). Therefore the optimal post-report duration is
\[
 T_*=
 \begin{cases}
 2\sqrt{r/a}, & r\le V^2/a,\\
 r/V+V/a, & r\ge V^2/a .
 \end{cases}
\]
Motion before the report cannot make the best freely chosen standby
placement faster.
\end{theorem}
\begin{proof}
Let \((x_0,v_0)\) be the common state when the report arrives and let
\(x_j(T)\) be a terminal position. Set \(\tau=\min(T/2,V/a)\).
Compare each trajectory with the same reference velocity, \(v_0\) up to
\(\tau\) and zero thereafter. Before \(\tau\), the difference is bounded
by \(at\), by the acceleration bound. After \(\tau\), the difference is
bounded by \(\min(a(T-t),V)\), by terminal rest and the speed cap.
These bounds equal \(\min(at,a(T-t),V)\) on their respective intervals.
The reference velocity is a device for the estimate and need not be a
feasible trajectory. Integrating gives
\[
 \|x_j(T)-(x_0+\tau v_0)\|
 \le \int_0^T\min(at,a(T-t),V)\,dt
 =D(T).
\]
Project \(x_0+\tau v_0\) onto \(C\). If its projection is \(p\), the
projection inequality \(\langle x_0+\tau v_0-p,z-p\rangle\le 0\) for
\(z\in C\) follows by minimizing distance along the segment from \(p\) to
\(z\). Expanding squared distances gives
\(\|p-z\|\le\|x_0+\tau v_0-z\|\). Apply it to each \(x_j(T)\). This
proves necessity without assuming the common state is at rest.

For sufficiency choose a minimizing standby point in \(C\); it exists
by compactness. Each \(F_j\) is compact, so choose a closest point
\(z_j\) and write the distance as \(d_j\). For \(T>0\), move on the
segment from standby to \(z_j\) at speed
\((d_j/D(T))\min(at,a(T-t),V)\), taking the constant path if \(d_j=0\).
The profile starts and ends at rest, has integral \(d_j\), and stays
within both bounds; the segment stays in \(C\). For \(T=0\), feasibility
is exactly the existence of a common point. Solving \(D(T)=r\)
gives the formula.
\end{proof}

\paragraph{How many worlds need to be checked?}
When the \(F_j\) are convex, each allowed standby set
\(\{x\in C:\operatorname{dist}(x,F_j)\le D(T)\}\) is convex.
By Helly's theorem, checking all simultaneous subfamilies of at most
\(d+1\) suffices; each subfamily still requires one common standby.
This bound cannot be reduced in general. Take the service sets to be
the individual vertices of a regular \(d\)-simplex within \(C\).
With side length \(s\), the smallest enclosing radius of a
\(k\)-vertex face is \(s\sqrt{(k-1)/(2k)}\). Its centroid is equidistant
from the \(k\) vertices, and their mean squared distance to any other
point increases by the squared distance to the centroid.
The displayed value follows by summing squared distances between
ordered distinct vertices. Choose \(D(T)\) between the
\(d\)-vertex and \((d+1)\)-vertex values.

\paragraph{Pinned geometry.}
Use the first acute triple (indices \(1,2,3\)) among the first five
initial positions in the published 3D scenario. The authored model
keeps the relay in a ball of radius \(1.65\) around recipient \(2\).
Each new service ball has radius \(R=0.8\). All three anchors lie in
\(C\). If \(b\) is their smallest enclosing radius, then
\(r=\max(0,b-R)\). The lower bound is the triangle inequality at every
anchor. A smallest-ball center lies in the convex hull of the anchors
(project onto their hull if necessary); the segment from it to each
anchor lies in \(C\) and reaches the corresponding service region
within distance \((b-R)_+\).

The replay certifies the enclosing ball by positive weights
\(\lambda_j\) whose sum is one. Their weighted centroid \(c\) is
equidistant from all three anchors \(p_j\), at squared distance \(b^2\).
For any \(x\), \(\sum_j\lambda_j\|x-p_j\|^2=b^2+\|x-c\|^2\).
The pairwise tests allow durations starting at approximately \(0.55105\)
with \(a=1,V=0.2\); the joint deadline is \(0.57816\). At duration
\(0.57\), every pair passes with a pair-specific standby and the three
together fail, even allowing pre-report motion. The decisions use
rational squared distances; the displayed cutoff times are rounded.

\paragraph{Boundaries.}
Welzl's enclosing-ball method is an applicable direct geometric
baseline for the equal-radius example; it gives the same placement and
radius. The contribution here is the exact reduction from causal motions
with a possibly nonzero common velocity, not a better enclosing-ball
algorithm. The source formation controller handles managed departures
with bridging edges formed in advance. Only its initial geometry is
reused; this is neither a simulation of that controller nor a ranking
against it. If \(C\) is not convex, the distance condition need not
suffice: take \(C=\{-1,1\}\) with a different terminal in each component.
If free standby placement is removed, the motion lower bound still
applies, but need not be attainable. This result does not discharge R6
or certify an independent mission.
\end{document}
